\documentclass[11pt]{article}
\usepackage[margin=2.6cm]{geometry}
\usepackage[T1]{fontenc}
\usepackage[utf8]{inputenc}
\usepackage{lmodern}
\usepackage{amsmath,amssymb}
\usepackage{array,longtable}
\usepackage{listings}
\usepackage{xcolor}
\usepackage[hidelinks]{hyperref}
\usepackage{parskip}

\lstset{
  basicstyle=\ttfamily\small,
  breaklines=true,
  columns=fullflexible,
  keepspaces=true,
  frame=single,
  framerule=0.2pt,
  xleftmargin=0.6em,
  aboveskip=0.7em,
  belowskip=0.7em
}
\lstdefinelanguage{RustLite}{
  morekeywords={let,mut,fn,pub,use,struct,enum,match,for,in,if,else,return},
  morecomment=[l]{//},
  morestring=[b]",
  sensitive=true
}

\newcommand{\code}[1]{\texttt{#1}}

\title{\code{physical\_object\_simulator}\\[2pt]
       \large The Complete Solution, Explained}
\author{physical\_object\_simulator}
\date{July 22, 2026}

\begin{document}
\maketitle
\tableofcontents

\emph{Written for a reader who has never seen this project, Rust, or a
numerical integrator before. The companion document \code{grammar.pdf}
specifies the command language; this one explains the whole system
underneath it.}

\section{What problem does this solve?}

The repository previously contained \textbf{three separate, overlapping
ways} to describe a physical thing:

\begin{center}
\begin{tabular}{llp{5.4cm}l}
\textbf{legacy type} & \textbf{file} & \textbf{what it modeled} & \textbf{its integrator}\\\hline
\code{PointParticle} & \code{src/solver.rs} & gravitating point mass; momentum, angular momentum, Laplace-vector observables & hand-written velocity Verlet\\
\code{RigidBody} & \code{RigidBody.rs} & full 3-D rigid body: quaternion, inertia tensors, charge, magnetic coupling, SDF boundary & hand-written semi-implicit Euler\\
\code{RigidBody3D} & \code{RigidBody3D.rs} & simpler rigid body: arrays, diagonal inertia from shape, sphere/cuboid enum & hand-written explicit Euler\\
\end{tabular}
\end{center}

The task: refactor into \textbf{one} pure-Rust simulator in
\code{./physical\_object\_simulator}, whose single object type
\code{physical\_object} is the \emph{unique union} of all three (every
data member and property, nothing duplicated), whose \emph{only}
numerical backend is the local \code{./sundials\_rs} library --- zero
\code{unsafe}, zero external dependencies, zero warnings --- driven by a
\emph{lexer + grammar + stack machine} command language and a
\emph{notebook} interface, bridged to \emph{JupyterLab}.

\subsection{Why ``sundials-only'' is a big deal}

SUNDIALS is Lawrence Livermore's industrial-strength suite of
differential-equation solvers. \code{../sundials\_rs} is a pure-Rust,
line-faithful translation of SUNDIALS 7.8.0 --- itself zero-unsafe and
dependency-free --- providing CVODE (adaptive multistep) and ARKODE
(including symplectic integrators). The legacy Euler/Verlet steppers
were \emph{not carried over}: every trajectory now comes from a solver
with proper error control or symplectic structure. The difference is not
academic --- explicit Euler gains energy every orbit; CVODE at default
tolerances conserves the outer solar system's energy to a part in
$10^{6}$ over 1{,}369 years (\S\ref{sec:verify}).

\section{The layout}

\begin{lstlisting}
rustSimulate/
|- Cargo.toml            workspace: two member crates, no external deps
|- PLAN.md               design record
|- grammar.md/.tex/.pdf  command-language spec
|- physical_object_simulator.md/.tex/.pdf   this document
|- ARCHITECTURE.md       module architecture & contracts
|- CLAUDE.md             contributor/agent working rules
|- physical_object/      THE LIBRARY
|    src/linalg.rs           Vec3, Mat3, Quat (pure std)
|    src/boundary.rs         Boundary enum + Sdf trait + inertia
|    src/physical_object.rs  pub struct physical_object (the union)
|    src/system.rs           PhysicalObjectSystem
|    src/integrate.rs        ALL integration (CVODE + ARKODE SPRK)
|    examples/ + tests/
|- posim/                THE FRONT END (binary posim)
|    src/{lexer,parser,vm,notebook,machine,main}.rs
|    src/scene/          THE GRAPHICAL SCENE WINDOW
|        mod.rs              server + playback thread + shared state
|        ws.rs               hand-rolled SHA-1 / base64 / RFC 6455
|        scene.html          the embedded window page (toolbar, canvas,
|                            status bar -- served to your browser)
|- jupyter/              JupyterLab wrapper kernel + tests
|- sundials_rs/          the pure-Rust SUNDIALS engine (own workspace)
\end{lstlisting}

Dependencies point one way:
\code{posim} $\to$ \code{physical\_object} $\to$
\{\code{sundials\_core}, \code{cvode\_rs}, \code{arkode\_rs}\}. The
workspace \code{Cargo.lock} lists \textbf{exactly five crates} --- the
machine-checkable proof of ``zero external dependencies.''

\section{The unique union: \code{pub struct physical\_object}}\label{sec:union}

``Unique union'' means: take every data member and property of the
three legacy types; where two describe \emph{the same physical
quantity}, store it \textbf{once} and reconcile the interfaces. Twelve
stored fields result:

\begin{longtable}{llp{8.2cm}}
\textbf{field} & \textbf{from} & \textbf{unification decision}\\\hline
\code{id} & PointParticle & user label\\
\code{mass} & all three & coupled to \code{inverse\_mass}\\
\code{inverse\_mass} & RigidBody3D & $m\le 0 \Rightarrow 1/m := 0$ = \textbf{static body}\\
\code{charge} & RigidBody(3D) & \\
\code{position} & all three & one \code{Vec3} replaces nalgebra and \code{[f64;3]}\\
\code{orientation} & RigidBody & always unit (setters renormalize)\\
\code{momentum} & RigidBody(3D) & \textbf{canonical}. PointParticle stored \code{velocity}; storing both would duplicate one quantity, so velocity became a \emph{derived property}: reads $p\,m^{-1}$, writes $p := m v$\\
\code{angular\_momentum} & RigidBody(3D) & the \textbf{spin} state; PointParticle's \emph{orbital} method collided with this name $\to$ renamed \code{orbital\_angular\_momentum(com)}; \code{total\_angular\_momentum} = orbital + spin\\
\code{inertia\_tensor} & RigidBody + 3D & one body-frame tensor\\
\code{inverse\_inertia\_tensor} & both & coupled; singular $\to$ zero inverse = ``cannot rotate''\\
\code{magnetic\_moment\_tensor} & RigidBody & torque $(RMR^{\mathsf T})B$\\
\code{boundary} & both & RigidBody had \code{Box<dyn Boundary>} (SDF trait); RigidBody3D an enum. The \textbf{enum} keeps the name --- \textbf{seven variants} now: the legacy \code{Point}/\code{Sphere}/\code{Cuboid} plus \code{Torus\{ring\_radius, tube\_radius\}}, \code{Disk\{radius\}}, \code{Cylinder\{radius, half\_height\}} and the compound \code{Dumbbell\{r1, r2, rod\_radius, z1, z2, f1, f2\}} (two solid spheres plus a solid rod as ONE rigid body; the stored mass fractions \code{f1}/\code{f2} keep every part mass recoverable), each with an exact SDF; the SDF behavior survives as the \code{Sdf} trait (verbatim central-difference normal) implemented \emph{for} the enum --- no trait objects, struct stays \code{Clone}\\
\end{longtable}

\textbf{Every field has public \code{get\_x()}/\code{set\_x()}} (coupled
ones maintain their invariants on every write), plus derived
\code{velocity} and \code{angular\_velocity}
($\omega = R\,I^{-1}R^{\mathsf T}L$) properties.

Methods carried over verbatim: \code{momentum()},
\code{orbital\_angular\_momentum(com)}, \code{laplace\_vector(com, k)}
(incl.\ the $r=0$ guard), \code{linear\_velocity()},
\code{angular\_velocity()}, \code{kinetic\_energy()}
($=\tfrac12 m|v|^2+\tfrac12\omega\cdot L$),
\code{to\_local\_space}/\code{to\_world\_space},
\code{signed\_distance}/\code{surface\_normal},
\code{recompute\_inertia\_from\_boundary()} (sphere $\tfrac25 mr^2$;
cuboid $\tfrac m3(h^2{+}h^2)$ diagonals; Point $\to$ zero; the shapes
added by the TORUS/DISK/CYLINDER and DUMBBELL releases get the closed
forms below).

Each of the non-legacy shapes has an exact SDF and an exact analytic
inertia tensor (body frame; $z$ is the shape's symmetry axis):

\begin{center}
\begin{tabular}{lp{4.6cm}p{6.4cm}}
\textbf{shape} & \textbf{parameters} & \textbf{analytic inertia}\\\hline
\code{Torus} & \code{ring\_radius} $c$ (centerline circle), \code{tube\_radius} $a$ & $I_z = m\,(c^2 + \tfrac34 a^2)$;\quad $I_x = I_y = m\,(\tfrac12 c^2 + \tfrac58 a^2)$\\
\code{Disk} & \code{radius} $a$ --- ideal, zero thickness & $I_x = I_y = \tfrac14 m a^2$;\quad $I_z = \tfrac12 m a^2$ (perpendicular-axis theorem: $I_z = I_x + I_y$)\\
\code{Cylinder} & \code{radius} $r$, \code{half\_height} $h$ (full height $2h$) & $I_z = \tfrac12 m r^2$;\quad $I_x = I_y = m\,(3r^2 + 4h^2)/12$\\
\code{Dumbbell} & sphere radii $r_1$/$r_2$ at $(0,0,z_1)$/$(0,0,z_2)$, \code{rod\_radius}, mass fractions $f_1$/$f_2$ & exact \textbf{composite}: $\tfrac25 m r^2$ for each sphere plus its parallel-axis term $m z^2$, plus the rod's cylinder terms about the COM\\
\end{tabular}
\end{center}

The dumbbell --- two solid spheres joined by a solid rod, \textbf{one}
rigid body --- is built by
\code{boundary::dumbbell(m1, m2, m\_rod, r1, r2, rod\_radius, length)},
which places the sphere offsets at
$z_1 = -(m_2 + m_{\mathrm{rod}}/2)\,L/M$ and
$z_2 = (m_1 + m_{\mathrm{rod}}/2)\,L/M$ ($L$ the center-to-center
length, $M = m_1 + m_2 + m_{\mathrm{rod}}$) so that the body-frame
origin \textbf{is} the center of mass: the identity
$m_1 z_1 + m_2 z_2 + m_{\mathrm{rod}}(z_1 + z_2)/2 = 0$ holds exactly
(pinned by test). Its SDF is the exact \textbf{union} --- the min of
the parts' SDFs --- and the stored mass fractions make every part mass
recoverable from the total, which is what lets the grammar's
\code{d.m1}/\code{.m2}/\code{.m\_rod}/\code{.r1}/\code{.r2}/%
\code{.rod\_radius}/\code{.length} member paths read \emph{and} write
the parts (a member write rebuilds total mass, COM offsets and the
inertia tensor in one step).

Every variant also answers an exact \textbf{support function} in
\code{boundary.rs}: \code{support\_extent(u)} $= \max_{x\in B}\,x\cdot u$
--- the farthest reach along a direction, a genuinely \textbf{directed}
quantity since the dumbbell release (the dumbbell is the first
non-centrally-symmetric shape: its off-center spheres make
$h(u) \ne h(-u)$; every symmetric shape's closed form is unchanged; for
the torus it is the support of its convex hull: a support function
cannot see the hole) ---
\code{support\_point(u)}, a point achieving that extent, which returns
the \textbf{centroid of the supporting set} when that set is a whole
face, edge, circle or cap (so a flat-on contact puts its contact point
at the face center and carries no spurious lever arm),
\code{support\_rank(u)} (the supporting set's dimension: 0 = vertex,
1 = edge/circle, 2 = face/cap), and \code{bounding\_radius()}. The
collision tiers of \S\ref{sec:collisions} and the anti-tunnel step cap
are built on exactly these.

Three constructors mirror the legacy \code{new} functions:
\code{new\_point}, \code{new\_rigid}, \code{new\_from\_shape}.

Naming quirk: the struct is deliberately \code{physical\_object}
(lower-case, per the specification; legal because crate roots allow
\code{non\_camel\_case\_types}) and lives in a same-named module, so the
import is
\code{use physical\_object::physical\_object::physical\_object;} and
sibling paths need a leading \code{::}.

\section{The system: \code{PhysicalObjectSystem}}

The legacy \code{GravitationalSystem}, generalized: objects,
\code{g\_constant}, \code{softening} (legacy const $\to$ field,
$10^{-6}$), uniform gravity/E/B fields, per-object external
force/torque, CVODE tolerances, \code{method}, \code{time}.
Observables ported verbatim: \code{total\_mass},
\code{center\_of\_mass}, \code{compute\_accelerations} (identical
softened arithmetic), plus \code{total\_momentum},
\code{total\_angular\_momentum}, \code{total\_energy},
\code{laplace\_vector(i)} ($k = G\,M_{\text{total}}$, about the center
of mass).

Solver packing --- \textbf{13 numbers per object}:
\begin{lstlisting}
[ x y z | px py pz | qw qx qy qz | Lx Ly Lz ]   per object -> length 13N
\end{lstlisting}
\code{pack\_state}/\code{unpack\_state} are exact inverses (unpack
renormalizes quaternions).

\section{Integration: how the physics advances}

\code{integrate.rs} is the \textbf{only} file that advances time. It
solves
\begin{align*}
\dot q &= p\,m^{-1}\\
\dot p &= \textstyle\sum_j G m_i m_j\,
          \frac{q_j-q_i}{(|\Delta|^2+\varepsilon^2)^{3/2}}
          + m\,g + qE + q\,v\times B + F_{\text{ext}}\\
\dot{\hat q} &= \tfrac12\,(0,\omega)\otimes\hat q,
   \qquad \omega = R\,I^{-1}R^{\mathsf T}L\\
\dot L &= \tau_{\text{ext}} + (R M R^{\mathsf T})B
\end{align*}

\subsection{CVODE (Adams or BDF) --- the general path}
\code{Method::Adams} (default) and \code{Method::Bdf} use CVODE with
Newton iteration and the dense linear solver (difference-quotient
Jacobian), following the reference pattern of
\code{cvode\_rs/examples/solar\_system.rs}. Contracts:
\begin{itemize}
\item \textbf{Callbacks are plain \code{fn} pointers} (no closures):
      parameters are cloned into a snapshot, passed as
      \code{user\_data: Option<Box<dyn Any>>}, downcast inside; a failed
      downcast returns $-1$ (SUNDIALS ``unrecoverable''), never panics.
\item \textbf{Quaternion drift}: renormalized only at output points and
      only when $\bigl||\hat q|-1\bigr| > 10^{-10}$; the state mutation
      invalidates CVODE's multistep history, so the driver then calls
      \code{CVodeReInit} (accumulating statistics first --- ReInit
      zeroes the counters).
\end{itemize}

\subsection{ARKODE SPRKStep --- the symplectic path}
\code{Method::Sprk\{table, dt\}} uses symplectic partitioned
Runge--Kutta at fixed step (state repacked \code{[q(3N) | p(3N)]},
following \code{ark\_kepler.rs}). Symplectic error stays \emph{bounded
forever} instead of drifting. The legacy velocity Verlet \emph{is}
\code{ARKODE\_SPRK\_LEAPFROG\_2\_2}. Separability is \textbf{gated}:
any $B$ field, magnetic tensor, external torque or spinning body yields
an error naming the offending feature.

\subsection{The API}
\begin{lstlisting}[language=RustLite]
pub fn run (system: &mut PhysicalObjectSystem, t_end: f64, nout: usize)
    -> Result<RunReport, String>;
pub fn step(system: &mut PhysicalObjectSystem, dt: f64)
    -> Result<RunReport, String>;
pub fn propagate_single(obj: &mut physical_object, force: Vec3,
    torque: Vec3, dt: f64) -> Result<(), String>;
\end{lstlisting}
\code{RunReport} carries per-output snapshots (time, energy, momentum,
angular momentum, center of mass) plus solver statistics
(\code{nst}, \code{nfe}, \code{nni}, \code{netf}).
\code{physical\_object::integrate(force, torque, dt)} --- the union of
the two legacy Euler methods --- delegates to
\code{propagate\_single}, so even ``one object, one step'' goes through
CVODE.

\section{The front end in brief}

(Fully specified in \code{grammar.pdf}.) \code{posim} runs the
lexer $\to$ parser $\to$ stack-machine pipeline; the notebook REPL
numbers cells \code{In[n]}/\code{Out[n]} with
\code{\%history}/\code{\%edit}/\code{\%rerun}/\code{\%save}/\code{\%load}
magics; \code{--script} replays files; \code{--notebook} replays a
file and then stays interactive (the dynamic-notebook launcher,
scene window alive); \code{--machine} speaks
line-delimited JSON (hand-rolled, zero dependencies). The JupyterLab
bridge is a $\sim$100-line Python wrapper kernel subprocessing
\code{posim --machine}.

\subsection{The graphical scene window}

\code{SCENE CREATE} opens a \textbf{separate graphical window, outside
the notebook} --- your web browser --- showing every simulator entity.
How can a zero-dependency program open a graphical window? It doesn't
link a GUI library at all: posim starts a tiny \textbf{HTTP +
WebSocket server} on \code{127.0.0.1} (written with nothing but
\code{std::net}, including the SHA-1 and base64 the WebSocket handshake
needs, in \code{posim/src/scene/ws.rs}) and serves one self-contained
HTML/canvas page (\code{scene.html}, embedded into the binary at
compile time). The browser supplies the pixels; posim supplies
everything else. Zero \code{unsafe}, zero external crates,
\code{Cargo.lock} still lists five.

The moving parts, and who talks to whom:

\begin{lstlisting}
 notebook / JupyterLab                     your web browser
   |  SCENE commands                        (the scene window)
   v                                             ^   | JSON text frames
  VM --> SceneHandle --> Arc<Mutex<Shared>> <----+   | over one WebSocket
                              ^                      v
                    playback thread (~30 fps):   toolbar clicks, camera
                    steps the scene's system     gestures, errors -- all
                    via physical_object::integrate   sent back as events
\end{lstlisting}

\begin{itemize}
\item \textbf{The playback thread} owns a \emph{synchronized copy} of
      the notebook's system and evolves it forward at
      \code{SCENE SET\_TIME\_STEP}'s \code{dt}, $\sim$30 frames per
      second, broadcasting each frame's positions and orientations to
      every connected window. Time is advanced \textbf{only} through
      \code{physical\_object::integrate} --- the sundials-only rule
      holds in the scene too; there is no second stepper hiding in the
      GUI.
\item \textbf{Reverse} (\code{SCENE REVERSE} or the toolbar button)
      replays a ring buffer of up to 20{,}000 snapshots recorded while
      stepping forward --- a faithful rewind of states the solver
      actually produced, never ``integration with negative dt''.
\item \textbf{Asynchronous, both directions.} Notebook $\to$ window:
      scene deltas, camera commands, playback state. Window $\to$
      notebook: errors, data requests and user actions are queued as
      \emph{events} you read with \code{SCENE EVENTS}; under
      \code{--machine}, each event is also pushed immediately as an
      unsolicited \code{\{"event": ...\}} JSON line. The Jupyter kernel
      runs a background reader thread that routes ordinary replies to
      the requesting cell and streams event lines into the notebook as
      \code{[scene] ...} output --- the window can talk to you
      \emph{between} cells.
\item \textbf{In the window}: a toolbar (Start, Pause, Stop, Reverse,
      single-step, dt entry, zoom, reset view, grid/trails/labels
      toggles, \code{?} help) and a status bar (connection dot, mode,
      simulated time, dt, total energy, body count, hidden count,
      history depth, camera yaw/pitch/distance, fps). Arrow keys
      translate the view, left-drag rotates, the mouse wheel or
      \code{+}/\code{-} zooms; \code{H} shows the full cheat-sheet.
\end{itemize}

\section{How we know it is right}\label{sec:verify}

\begin{center}
\begin{tabular}{p{6.2cm}p{8.2cm}}
\textbf{check} & \textbf{result}\\\hline
\code{cargo build --workspace} & zero warnings (\code{\#![deny(warnings)]})\\
\code{cargo test --workspace} & 52 tests green (17 library, 9 conservation, 26 posim)\\
\code{\#![forbid(unsafe\_code)]} & every crate root; \code{grep unsafe} finds nothing else\\
\code{Cargo.lock} & exactly 5 local crates $\Rightarrow$ zero external deps\\
outer solar system & matches the donor sundials example: Pluto to \textbf{8 decimals after 500{,}000 days}; energy drift $7.8\times10^{-7}$\\
Kepler $e=0.6$ & energy, $L$, Laplace vector conserved $<10^{-6}$ on both paths\\
torque-free tumbling & $\Delta L = 0$ exactly; KE drift $5\times10^{-9}$; $|\hat q|$ drift 0\\
charged gyration & gyroradius $= mv/(qB)$ to $10^{-4}$; orbit closes to $7\times10^{-9}$\\
Jupyter & protocol test (24 checks incl.\ the SCENE family + \code{events} op), ZMQ kernel test (5 cells), live JupyterLab launch + cell execution\\
WebSocket layer & SHA-1 against the FIPS 180-4 test vectors; base64 against RFC 4648 \S10; the handshake against the worked example in RFC 6455 \S1.3\\
scene server & three integration tests drive a real TCP client: page serves with toolbar/status bar; a WebSocket session handshakes and streams frames; camera sync / start / pause / events round-trip; forward-then-reverse playback lands \textbf{exactly} back on the $t=0$ state\\
scene window (real browser) & 16 headless-Chrome checks dispatching genuine input events: arrow keys translate the view (and return exactly), left-drag changes yaw \emph{and} pitch (and does not zoom), wheel and \code{+}/\code{-} zoom in/out, toolbar Start evolves $t$, Pause freezes, Reverse plays $t$ backward; the status bar reports it all\\
\end{tabular}
\end{center}

\section{Fourteen more worked examples}

Additional to the command-language examples in \code{grammar.pdf}
--- these exercise the \textbf{Rust API}, the \textbf{machine protocol},
\textbf{JupyterLab}, and the \textbf{graphical scene window}. Every number was produced by the real code. To
run a Rust snippet, drop it into
\code{physical\_object/examples/demo.rs} (with \code{fn main()}) and
\code{cargo run -p physical\_object --example demo}.

Common imports:
\begin{lstlisting}[language=RustLite]
use ::physical_object::boundary::Boundary;
use ::physical_object::integrate::{run, step, Method};
use ::physical_object::linalg::{Quat, Vec3};
use ::physical_object::physical_object::physical_object;
use ::physical_object::{PhysicalObjectSystem, Sdf};
\end{lstlisting}

\subsection{S1 --- build, run, and read a report (the canonical loop)}

\begin{lstlisting}[language=RustLite]
let sun   = physical_object::new_point(0, 1.0e9, Vec3::zeros(), Vec3::zeros());
let earth = physical_object::new_point(1, 1.0,
    Vec3::new(1.0, 0.0, 0.0), Vec3::new(0.0, 1.0, 0.0));

let mut sys = PhysicalObjectSystem::new(vec![sun, earth], 1.0e-9); // G*M ~= 1
sys.softening = 0.0;
sys.method = Method::Adams;

let e0 = sys.total_energy();
let report = run(&mut sys, 6.28, 4)?;          // one orbit, 4 snapshots

for s in &report.snapshots {
    println!("t = {:5.2}  E-drift = {:.2e}  com = {:?}",
             s.t, ((s.energy - e0) / e0).abs(), s.center_of_mass.to_array());
}
println!("solver steps = {}, rhs evals = {}", report.nst, report.nfe);
\end{lstlisting}

\emph{What to notice.} \code{run} takes an \textbf{absolute} end time
and mutates the system in place; \code{RunReport} gives
conserved-quantity snapshots with no extra bookkeeping. All solver
errors are \code{Result<\_, String>} --- no panics in library code.

\subsection{S2 --- reading eccentricity off the Laplace vector}

For a unit-mass orbiter with $GM = 1$, the eccentricity \emph{vector}
is $\vec e = \vec A/(mk)$: its length is the orbital eccentricity and
it points at perihelion.

\begin{lstlisting}[language=RustLite]
let ecc = 0.6;
let central = physical_object::new_point(0, 1.0e9, Vec3::zeros(), Vec3::zeros());
let orbiter = physical_object::new_point(1, 1.0,
    Vec3::new(1.0 - ecc, 0.0, 0.0),                       // perihelion
    Vec3::new(0.0, ((1.0 + ecc) / (1.0 - ecc)).sqrt(), 0.0));
let mut s = PhysicalObjectSystem::new(vec![central, orbiter], 1.0 / (1.0e9 + 1.0));
s.softening = 0.0;

run(&mut s, 37.0, 5)?;                                    // ~6 orbits

let k = s.g_constant * s.total_mass();                    // = 1
let m = s.objects[1].get_mass();
let e_vec = s.laplace_vector(1).unwrap() / (m * k);
// prints: e_vec = [0.600000, 0.000000, 0.000000], |e| = 0.600000
\end{lstlisting}

\emph{What to notice.} After six orbits the recovered eccentricity is
\code{0.600000} --- six exact digits, using the \emph{legacy
PointParticle formula} ($A = p\times L - mk\hat r$,
$k = G M_{\text{total}}$), now demonstrably correct because the
integrator underneath no longer lies. The $k$-scaling makes the formula
the true conserved LRL vector only for a unit-mass orbiter --- hence
$m=1$ with $G$ scaled down.

\subsection{S3 --- retiring the legacy Euler: \code{integrate()} on one body}

The legacy \code{RigidBody3D} demo dropped a charged sphere under
gravity with a ``wind'' torque via explicit Euler at $dt=0.1$. The same
call now runs CVODE:

\begin{lstlisting}[language=RustLite]
let mut body = physical_object::new_from_shape(
    0, 2.0, -1.5,
    Vec3::new(0.0, 10.0, 0.0),        // position
    Vec3::new(1.0, 0.0, -0.5),        // velocity
    Vec3::new(0.0, 2.0, 0.0),         // angular velocity
    Boundary::Sphere { radius: 0.5 },
);
let gravity = Vec3::new(0.0, -9.81 * 2.0, 0.0);   // F = m g
let wind    = Vec3::new(0.1, 0.0, 0.0);

body.integrate(&gravity, &wind, 0.3)?;   // same signature as the legacy code
// exact: y(0.3) = 10 - 9.81*0.09/2 = 9.55855;  L = [0.03, 0.4, 0]
\end{lstlisting}

\emph{What to notice.} The method signature is the legacy one, so old
call sites read the same --- but the result is exact to solver
tolerance (the committed test asserts agreement to $10^{-8}$). At
$dt=0.3$, explicit Euler's position error would be
$\sim4\times10^{-2}$ --- six orders worse.

\subsection{S4 --- a charge in a uniform E field (exact answer)}

Force $qE$ on a resting charge gives $x(t) = \frac{qE}{2m}t^2$.

\begin{lstlisting}[language=RustLite]
let mut ion = physical_object::new_point(0, 2.0, Vec3::zeros(), Vec3::zeros());
ion.set_charge(3.0);
let mut s = PhysicalObjectSystem::new(vec![ion], 0.0);   // gravity off (G = 0)
s.e_field = Vec3::new(0.5, 0.0, 0.0);
run(&mut s, 4.0, 1)?;
// analytic: 3*0.5/(2*2) * 16 = 6
// prints:   x = 6.000000000000, analytic = 6.000000000000
\end{lstlisting}

\emph{What to notice.} Twelve matching digits: for polynomial solutions
a multistep method of sufficient order is \emph{exact} up to round-off.
Note $G=0$: every interaction (pairwise gravity, uniform fields,
Lorentz, external) is independently switchable.

\subsection{S5 --- static bodies: the \code{inverse\_mass = 0} convention}

\begin{lstlisting}[language=RustLite]
let mut anchor = physical_object::new_point(0, 5.0,
    Vec3::new(0.0, 10.0, 0.0), Vec3::zeros());
anchor.set_inverse_mass(0.0);          // static; get_mass() reports 0
let mut s = PhysicalObjectSystem::new(vec![anchor], 0.0);
s.uniform_gravity = Vec3::new(0.0, -9.81, 0.0);
run(&mut s, 3.0, 1)?;
// prints: [0.0, 10.0, 0.0]  -- three seconds of gravity, zero motion
\end{lstlisting}

\emph{What to notice.} The coupled setters keep the pair consistent
both ways: \code{set\_inverse\_mass(0)} back-computes $m=0$, so both
the force $mg$ and the conversion $p\,m^{-1}$ vanish --- pinned.
\textbf{Caveat:} with $m=0$ the body also stops \emph{sourcing}
pairwise gravity; use static bodies as anchors and test charges, not as
fixed suns (for a fixed sun use a huge mass ratio, as in S1/S2).

\subsection{S6 --- geometry: signed distances, normals, frames}

\begin{lstlisting}[language=RustLite]
let mut brick = physical_object::new_from_shape(0, 1.0, 0.0,
    Vec3::new(5.0, 0.0, 0.0), Vec3::zeros(), Vec3::zeros(),
    Boundary::Cuboid { half_extents: [1.0, 2.0, 3.0] });
brick.set_orientation(Quat::from_axis_angle(Vec3::new(0.0, 0.0, 1.0),
                                            std::f64::consts::FRAC_PI_2));

let world_p = Vec3::new(7.5, 0.0, 0.0);
let local_p = brick.to_local_space(&world_p);    // -> [0, -2.5, 0]
let sd      = brick.signed_distance(&local_p);   // -> 0.5
let n       = brick.surface_normal(&Vec3::new(0.0, -2.5, 0.0)); // -> [0,-1,0]
\end{lstlisting}

\emph{What to notice.} The brick is rotated $90^\circ$ about $z$, so a
point ``to its right'' in world space is ``below it'' in body space:
\code{to\_local\_space} performs $q^{-1}(p - x)$ and returns
$[0,-2.5,0]$. The local $y$ half-extent is 2, so the point sits 0.5
\emph{outside}: \code{signed\_distance} $= 0.5$ (verified). The normal
comes from the legacy central-difference default ($\varepsilon=10^{-6}$)
--- the RigidBody SDF machinery, alive inside an enum.

\subsection{S7 --- the SPRK separability gate as an API}

\begin{lstlisting}[language=RustLite]
let mut s = /* Kepler system as in S2 */;
s.b_field = Vec3::new(0.0, 0.0, 1.0);            // breaks separability
s.method = Method::Sprk { table: "ARKODE_SPRK_LEAPFROG_2_2".into(), dt: 1e-3 };

match run(&mut s, 10.0, 5) {
    Err(msg) if msg.contains("separable") => {
        eprintln!("SPRK refused: {msg}");
        s.method = Method::Bdf;                   // graceful fallback
        run(&mut s, 10.0, 5)?;
    }
    other => { other?; }
}
\end{lstlisting}

\emph{What to notice.} The error is a \emph{diagnosis}: ``SPRK method
requires a separable Hamiltonian: magnetic field B must be zero (the
Lorentz force q v x B is velocity-dependent); use METHOD ADAMS or
BDF.'' Library code can pattern-match it; notebook users read the same
words. Nothing silently produces a wrong symplectic answer.

\subsection{S8 --- driving the simulator from Python (no Jupyter)}

\begin{lstlisting}
import json, subprocess
p = subprocess.Popen(["target/release/posim", "--machine"],
                     stdin=subprocess.PIPE, stdout=subprocess.PIPE,
                     text=True, bufsize=1)
def rpc(**req):
    p.stdin.write(json.dumps(req) + "\n"); p.stdin.flush()
    return json.loads(p.stdout.readline())

rpc(op="exec", code="new sphere { mass = 2, radius = 0.5, charge = -1.5 }")
rpc(op="set",  path="obj0.velocity", value=[1, 0, -0.5])
print(rpc(op="get", path="obj0.momentum")["result"])   # [2.0, 0.0, -1.0]
rpc(op="set",  path="system.uniform_gravity", value=[0, -9.81, 0])
rpc(op="exec", code="step 1")
state = rpc(op="state")["result"]     # whole system as one JSON document
\end{lstlisting}

\emph{What to notice.} \code{get}/\code{set} are direct wire versions
of the get/set API --- a scripting language away from batch parameter
sweeps; \code{state} dumps the entire system for logging or plotting
from any language. This exact flow is what
\code{jupyter/test\_protocol.py} asserts (11 checks, all green).

\subsection{S9 --- the same physics from a JupyterLab notebook}

With the kernel installed (\code{jupyter/README.md}), each cell is
posim language; shift-enter executes. A real 3-cell session, as run
through the Jupyter messaging protocol during verification:

\begin{lstlisting}
Cell 1:  new sphere { mass = 2, radius = 0.5 }
         -> obj0
Cell 2:  set system.gravity = [0, -9.81, 0]
         set obj0.position = [0, 10, 0]
         step 1
         -> t = 1 (advanced by 1, 12 solver steps)
Cell 3:  get obj0.position.y
         -> 5.095000000000006
\end{lstlisting}

\emph{What to notice.} Cell 2 is \textbf{multi-line}: the kernel
forwards a cell one line at a time, stopping at the first error ---
errors arrive as red stderr text (the verification suite deliberately
executes \code{get obj0.bogus\_field} and checks that ``unknown object
field'' comes back). Kernel restart = fresh simulator; the notebook
file = your reproducible experiment.

\subsection{S10 --- choosing an integrator: an honest benchmark}

Same Kepler orbit ($\sim$80 revolutions to $t=500$), measured:

\begin{center}
\begin{tabular}{lrr}
\textbf{method} & \textbf{internal steps} & $|\Delta E/E|$ at $t{=}500$\\\hline
CVODE Adams (rtol $10^{-10}$) & 22{,}194 & $2.5\times10^{-8}$\\
SPRK leapfrog, $dt=0.01$ & 50{,}000 & $1.3\times10^{-4}$\\
SPRK McLachlan-4-4, $dt=0.001$ (2 orbits) & 12{,}600 & $9.2\times10^{-14}$\\
\end{tabular}
\end{center}

\emph{What to notice.} Over this horizon the high-order adaptive Adams
beats 2nd-order leapfrog on raw accuracy --- the symplectic advantage
is not ``smaller error'' but \textbf{bounded} error: leapfrog's
$1.3\times10^{-4}$ oscillates forever and never grows, while any
non-symplectic method's error drifts secularly. For $10^{6}$-orbit runs
symplectic wins; for a thousand orbits at tight tolerance Adams is
cheaper; McLachlan-4-4 is 4th-order \emph{and} symplectic. Rules of
thumb: \textbf{Adams} for smooth accuracy, \textbf{BDF} for stiffness
(gyration), \textbf{SPRK} for astronomical-length point-mass runs.

\subsection{S11 --- watching a tumbling cuboid in the scene window}

Everything here is typed into the ordinary notebook (\code{cargo run});
the \emph{graphics} happen in the browser window that opens.

\begin{lstlisting}
In[1]: new cuboid { mass = 2, half_extents = [0.3, 0.2, 0.1],
                    position = [0, 2.5, 0.5], angular_velocity = [3, 0.2, 0.1] }
Out[1]: obj0
In[2]: new sphere { mass = 256, radius = 0.35 }
Out[2]: obj1
In[3]: scene create
Out[3]: scene window created: http://127.0.0.1:41372/
        (opened in your browser; if no window appeared, open that
        address yourself)
        showing 2 entities; SCENE START begins the evolution
In[4]: scene set_time_step 0.002
Out[4]: scene time step dt = 0.002
In[5]: scene start
Out[5]: scene playback: running
\end{lstlisting}

\emph{What you see.} A dark 3-D viewport with a ground grid and the
world axes: \code{obj1} is a shaded sphere, \code{obj0} a
\textbf{wireframe box} that tumbles --- its long axis wobbles because
$\omega$ is not aligned with a principal axis (torque-free precession,
the same physics the \code{tumbling\_body} example checks numerically).
Colored \textbf{trails} show where each body has been; the status bar
counts simulated time up at $dt = 0.002$ per solver step, $\sim$30
steps a second. Now steer it:

\begin{lstlisting}
In[6]: scene hide 1              # sphere vanishes; the box stays
Out[6]: 1 object(s) hidden
In[7]: scene show all
Out[7]: 0 object(s) hidden
In[8]: scene rotate 45 -10       # orbit the camera: yaw +45, pitch -10
Out[8]: camera yaw = -15 deg, pitch = 45 deg
In[9]: scene zoom in
Out[9]: camera distance = 9.6
In[10]: scene pause
Out[10]: scene playback: paused
In[11]: scene reverse
Out[11]: scene playback: reversing
\end{lstlisting}

\emph{What to notice.} At \code{In[11]} the window plays the motion
\textbf{backward in time} --- the box un-tumbles along the exact
recorded states (snapshot replay, not re-integration), pausing by
itself when it reaches the beginning and reporting it:
\code{scene events} then prints \code{reverse: reached the beginning of
recorded history --- paused}. The same controls live in the window's
toolbar, and the keyboard/mouse work directly: arrow keys translate,
left-drag rotates, wheel zooms.

\subsection{S12 --- driving the scene from JupyterLab (async events)}

The scene window can talk to the notebook \emph{between} cells --- the
part a synchronous REPL cannot do. With the kernel installed:

\begin{lstlisting}
Cell 1:  new point { mass = 1, position = [1, 0, 0], velocity = [0, 1, 0] }
         new point { mass = 1000, position = [0, 0, 0] }
         scene create
         -> obj0
         -> obj1
         -> scene window created: http://127.0.0.1:35519/ ...

Cell 2:  scene start
         -> scene playback: running

   (you click Pause and then Stop in the browser window's toolbar;
    with no cell running, these lines appear in the notebook by
    themselves)

         [scene] window action: pause
         [scene] window action: stop

Cell 3:  scene status
         -> scene: http://127.0.0.1:35519/  (1 window(s) connected)
           mode = stopped, t = 3.762, dt = 0.01, steps = 376,
           history = 0 frame(s)
           entities = 2 (hidden: none)
           camera: yaw = -60, pitch = 55, dist = 12, target = [0, 0, 0]
\end{lstlisting}

\emph{What to notice.} The \code{[scene]} lines are
\textbf{asynchronous}: posim pushed
\code{\{"event":"scene","message":"window action: pause"\}} on stdout
the moment you clicked, the kernel's background reader thread caught
it, and JupyterLab printed it --- no cell was executing. Errors travel
the same road, marked red: if the window's JavaScript throws, or you
type a bad dt into its toolbar, the notebook shows e.g.\
\code{[scene] error: window sent a non-positive dt --- ignored} on
stderr. In the plain terminal notebook the same events queue up
silently instead (no async printing into your prompt) and
\code{scene events} drains them on demand --- one mechanism, two
delivery styles.

\subsection{S13 --- the box of shapes (every variant in one rigid box)}

The committed script
\code{scripts/collisions/11\_box\_of\_shapes.posim} rattles
\textbf{all six shapes} inside a \code{BOX 4}. At the center sits a
mass-1 torus (inner radius 1, outer radius 2) tilted onto the axis
$(1,1,1)/\sqrt3$ --- an \emph{axis-aligned} torus of outer radius 2
would exactly inscribe the 4-box, but tilted, its support extent per
world axis is $1.5\sqrt{2/3} + 0.5 \approx 1.7247$, clearance $0.2753$
(\S\ref{sec:union}'s \code{support\_extent}, and a committed test).
Around it: a point particle --- the \textbf{only mover},
$v = (100, 200, 100)$ --- a mass-2 sphere, a $2/3$-mass disk, a
$5/3$-mass cube and a mass-2 cylinder, positions drawn by a small
documented LCG inside the script. Captured session (abridged):

\begin{lstlisting}
In[2]:= box 4
Out[2]= box: inner size 4 x 4 x 4 --- six static walls obj0, obj1,
        obj2, obj3, obj4, obj5 with inverse_mass = 0 (infinitely
        massive); objects collide elastically off the inside faces
In[4]:= new point { mass = 1, position = [1.406590, -0.995859,
                    0.569601], velocity = [100, 200, 100] }
Out[4]= obj7
In[10]:= collide
Out[10]= collisions ON (51 collidable pair(s); 0 impulse(s) so far)
In[13]:= energy
Out[13]= 30000
In[14]:= momentum
Out[14]= [100, 200, 100]
In[15]:= run 0.1 steps 100
Out[15]= t = 0.1 (2121 solver steps, 100 snapshots, |dE/E| = 2.040e-10,
         119 collision(s) --- CONTACTS lists them)
In[16]:= energy
Out[16]= 29999.999993880127
In[17]:= momentum
Out[17]= [146.48911126803657, 102.1478121131382, 46.01601291636828]
\end{lstlisting}

\emph{What to notice.} $E_0 = \tfrac12\cdot 1\cdot(100^2 + 200^2 +
100^2) = \textbf{30000}$ exactly (one mover). In 0.1\,s the box
produces \textbf{119 collisions} spanning all three dispatch tiers of
\S\ref{sec:collisions}, and total energy survives every one of them to
$|\Delta E/E| = 2.040\times10^{-10}$. Momentum does \textbf{not}
survive --- $(100, 200, 100) \to (146.49, 102.15, 46.02)$ --- and that
is the \emph{physical signature} of the infinitely massive walls: they
absorb momentum without moving ($\Delta v = \Delta p\,m^{-1} =
\Delta p\cdot 0$), and after the run all six slabs are still
bit-identically at rest. The 51 collidable pairs are
$\binom{12}{2} = 66$ minus the 15 static wall--wall pairs, which are
skipped. In the scene window the same setup shows the new wireframes
--- torus, disk, cylinder, quaternion-rotated --- plus a dashed
interior box outline; the wall slabs themselves are not drawn as
bodies.

\subsection{S14 --- two dumbbells: a user-defined function builds them,
and $E$, $P$ and $L$ survive the impact}

The committed script \code{scripts/collisions/12\_two\_dumbbells.posim}
defines \code{create\_dumbell(...)} \textbf{in the notebook language}
(\code{DEF} --- eleven parameters, all but the name with defaults) and
calls it twice to build two \textbf{named} rigid dumbbells --- two solid
spheres joined by a solid rod, one rigid body whose local origin is its
center of mass (\S\ref{sec:union}). They approach off-center with spin
and collide twice, elastically, with $G = 0$. Abridged captured session
(the full transcript is Example 12 of \code{collision\_detection.pdf}):

\begin{lstlisting}
In[3]:= create_dumbell("dumbell0", 1, 2, 0.5, 0.25, 0.25, 0.1, 1,
                       [-2, 0.15, 0], [1.5, 0, 0], [0, 0, 0.6])
Out[3]= obj0 as dumbell0
In[4]:= create_dumbell("dumbell1", 2, 1, 0.4, 0.3, 0.2, 0.08, 1.2,
                       [2, -0.15, 0], [-1.5, 0, 0], [0.4, 0, 0])
Out[4]= obj1 as dumbell1
In[9]:= energy
Out[9]= 7.865310611764706
In[10]:= momentum
Out[10]= [0.15000000000000036, 0, 0]
In[11]:= angmom
Out[11]= [0.4443030588235295, 0, -1.5059999999999998]
In[12]:= run 3 steps 60
Out[12]= t = 3 (3861 solver steps, 60 snapshots, |dE/E| = 9.463e-11,
         2 collision(s) --- CONTACTS lists them)
In[13]:= energy
Out[13]= 7.865310611020375
In[14]:= momentum
Out[14]= [0.15000000000000036, 0, 0]
In[15]:= angmom
Out[15]= [0.44430305882373156, -0.000000000000009547918011776346,
         -1.505999999999855]
\end{lstlisting}

\emph{What to notice.} Through two real CVODE collision events, energy
drifts by $|\Delta E/E| = 9.463\times10^{-11}$ (printed by the run line
itself), momentum is \textbf{bit-identical}, and total angular momentum
about the origin drifts at the $10^{-13}$ level --- \textbf{all three
conserved at once}, unlike Example S13, where the infinitely massive
walls legitimately ate momentum. Nothing here is static, and each
impulse pair $\pm J\hat n$ acts at \textbf{one shared contact point}, so
the two angular impulses about any origin cancel exactly:
action--reaction produces zero net torque (the anchor test
\code{colliding\_dumbbells\_conserve\_energy\_momentum\_and\_angular\_momentum}
pins $E$, $P$ and $L$ to $10^{-8}$). In the scene window the entity
labels show the registered user names (\code{dumbell0},
\code{dumbell1}), the dumbbells draw as two shaded spheres at their
rotated COM offsets joined by the rod's four silhouette lines, and the
window's permanent labeled ``conserved quantities'' readout displayed
\code{E = 7.86531061, P = [0.15000, 0.00000, 0.00000] |.| = 0.15000,
L = [0.44430, 0.00000, -1.50600] |.| = 1.57017} identically before and
after the impact (after it, $L$'s $y$ component read
$-1.31228\times10^{-13}$).

\section{Building, running, testing}

\begin{lstlisting}
cd rustSimulate
cargo run                              # the notebook (type HELP)
cargo test --workspace                 # all 104 tests (incl. scene server)
cargo build --release                  # -> target/release/posim
cargo run -p physical_object --release --example kepler_orbit
cargo run -p physical_object --release --example outer_solar_system
cargo run -p physical_object --release --example tumbling_body
cargo run -p physical_object --release --example charged_in_b_field
cargo run -p posim -- --script my_session.posim
cargo run -p posim --release -- --notebook dynamic_notebooks/kepler_orbit.posim
cargo run -p posim -- --machine

# the graphical scene window: inside the notebook type
#   scene create        (opens http://127.0.0.1:<port>/ in your browser)
# POSIM_NO_BROWSER=1 suppresses the browser launch (headless runs/CI);
# the URL is always printed, so you can open it by hand.

python3 jupyter/test_protocol.py       # machine-protocol checks (stdlib
                                       # only, incl. SCENE + events op)
\end{lstlisting}

\section{Collisions (new)}\label{sec:collisions}

Rigid-body collision detection is \textbf{on by default in every
scene}: \code{STEP}, \code{RUN} and the scene window detect impacts
\emph{during} the time step by SUNDIALS event rootfinding (the
integrator lands on the exact time of impact), resolve them with an
impulse along the \textbf{contact normal} --- the action--reaction
line, exposed as \code{contactK.normal} --- and continue. Per-object
\code{restitution} (default 1 = elastic), the
\code{COLLIDE}/\code{CONTACTS} commands, and the full science
reference (seven simulators surveyed, comparison chart, porting
evaluation, twelve worked examples with captured output) live in
\code{collision\_detection.pdf}. Quick taste:

\begin{lstlisting}
In [1]: new sphere { mass = 1, radius = 0.5, position = [-2,0,0], velocity = [1,0,0] }
In [2]: new sphere { mass = 1, radius = 0.5, position = [2,0,0], velocity = [-1,0,0] }
In [3]: step 3
Out[3]= t = 3 (advanced by 3, 26 solver steps, 1 collision(s) --- CONTACTS lists them)
In [4]: get contact0.normal
Out[4]= [1, 0, 0]
\end{lstlisting}

\textbf{The shape tiers.} With the \code{Boundary} enum at seven
variants (\S\ref{sec:union}; created with
\code{NEW TORUS|DISK|CYLINDER|DUMBBELL \{ ... \}} --- \code{grammar.pdf}
has the parameter grammar), \code{collide.rs} dispatches every pair in
three exactness tiers --- still no GJK/EPA machinery anywhere:

\begin{enumerate}
\item \textbf{Ball vs anything --- exact, via SDF closest points.} A
      \emph{ball} (a \code{Sphere}, or a \code{Point} as the
      zero-radius case) is tested against every shape through that
      shape's exact signed-distance field: the separation is
      $\mathrm{sdf}(\text{center}) - r$, and the contact normal and
      point come from the exact closest surface point. Because the
      true distance field is used, a small ball genuinely
      \textbf{threads a torus hole}, and a point particle passes
      straight through the ideal zero-thickness disk (its unsigned
      distance never crosses zero --- see the Limitations section).
\item \textbf{Cuboid vs cuboid --- exact 15-axis SAT}, unchanged from
      the collision release.
\item \textbf{Extended vs extended --- support-axis tests.} Every
      remaining pair (torus/disk/cylinder against each other or
      against a cuboid) is tested along candidate axes: each cuboid's
      three face axes, each round shape's symmetry axis, their
      pairwise cross products, the radial rejection axes (the
      component of the center offset perpendicular to each primary
      axis --- the true lateral direction when round shapes sit side
      by side with parallel axes), and the center line, with the gap
      computed from \S\ref{sec:union}'s exact support extents. Exact
      for face-on contacts --- in particular \textbf{every wall-face
      contact} --- and exact for \textbf{side-side contacts of
      parallel and near-parallel round shapes} (two parallel cylinders
      separate by exactly their lateral gap); only a genuinely
      \textbf{skew corner-on} approach remains conservative (contact
      may register slightly early, along the nearest candidate axis).
      The contact point is the \textbf{support point of the
      lower-support-rank body} --- a tilted cylinder against a wall
      face contributes its rim point, not the wall's face centroid;
      when the ranks tie, the body with the clearly smaller flat
      footprint wins, so a small cap landing on a big wall face
      contacts at the cap center. At this tier the torus is its convex
      hull: only balls thread the hole.
\end{enumerate}

\textbf{Compound bodies decompose.} The dumbbell never reaches a tier
as a whole: the narrow phase decomposes dumbbell-vs-anything over its
parts --- each sphere part is a ball tested through the other shape's
exact SDF (tier 1, exact against everything, \emph{including another
dumbbell}, whose union SDF is exact), the rod recurses as a
free-standing cylinder, and the deepest part supplies the contact; only
\textbf{rod-vs-rod} ever reaches the approximate tier 3. Dynamically
the dumbbell stays \textbf{one} rigid body, and because the impulse
pair acts at one shared contact point, two colliding dumbbells conserve
$E$, $P$ \emph{and} $L$ (Example S14).

\textbf{The rigid box: \code{BOX <size>} $|$ \code{BOX OFF} $|$
\code{BOX}.} One command walls the world in: six \textbf{static
\code{Cuboid} wall slabs} enclosing an inner
$\mathit{size}\times\mathit{size}\times\mathit{size}$ cube.
``Infinitely massive'' is exact, not an approximation, because the
equations of motion only ever consume the \textbf{inverse} mass:
velocity is read as $v = p\,m^{-1}$ (\S\ref{sec:union}), and the
collision impulse divides by
$n\!\cdot\!Kn = m_i^{-1} + m_j^{-1} + (\text{angular terms})$. A wall
has \code{inverse\_mass = 0} and zero inverse inertia, so it
contributes \textbf{0} to that denominator and receives no state
writes --- bodies bounce elastically off the inside faces while the
walls stay bit-identically at rest. The measurable consequence: energy
is conserved, \textbf{momentum is not} --- the walls absorb it
($\Delta p$ is finite but $\Delta v = \Delta p\cdot 0 = 0$), exactly
as a rigid room should. Bare \code{BOX} reports status,
\code{GET system.box} reads the inner size (0 = none), and \code{LIST}
tags each slab \code{[wall: static, inverse\_mass=0]}. The slabs are
ordinary objects with ordinary \code{objN} handles (see the
Limitations section). Example S13 puts all six shapes inside one.

\section{Limitations (told straight)}

\begin{itemize}
\item \textbf{Collision response is restitution-only} (normal
      impulses; \S\ref{sec:collisions}). Coulomb friction is designed
      and documented in \code{collision\_detection.pdf} but not
      shipped yet; contacts are single deepest-point (no persistent
      multi-point manifolds), so tall box stacks are not this build's
      target.
\item \textbf{Support-axis contacts are conservative on skew corners.}
      Extended-vs-extended pairs (torus/disk/cylinder against each
      other or a cuboid) are tested along candidate axes only
      (\S\ref{sec:collisions}, tier 3): exact for face-on and
      wall-face contacts \textbf{and} for side-side contacts of
      parallel/near-parallel round shapes (the radial rejection axes),
      but a genuinely skew corner-on approach may register contact
      slightly early, along the nearest candidate axis --- a
      convex-hull-level answer. In particular that tier sees the
      torus's convex hull, so \textbf{only balls (spheres/points) can
      thread the torus hole}. The dumbbell decomposes over its parts
      (\S\ref{sec:collisions}), so its sphere contacts are exact
      against everything; only \textbf{rod-vs-rod} --- two dumbbells
      meeting rod against rod --- remains a support-axis pair and
      shares the skew-corner caveat.
\item \textbf{The ideal disk is transparent to point particles --- and
      to a parallel ideal disk.} A zero-thickness disk meets a
      zero-radius point in a measure-zero event --- the disk's
      unsigned distance touches 0 without crossing it, so the
      rootfinder sees no sign change and the point passes through. Two
      disks with \textbf{parallel planes} are the same case squared:
      both have zero extent along the shared normal, so their
      separation is $|dz|$ and a face-on disk-disk crossing produces
      no root either (pinned by the test
      \code{parallel\_disk\_disk\_separation\_is\_the\_documented\_limitation}).
      Any ball of radius $r > 0$ collides with a disk normally; give a
      ``point'' a tiny sphere radius --- or tilt one disk, or model it
      as a thin cylinder --- if you need the crossing to fire.
\item \textbf{\code{BOX} walls are real objects.} The six slabs occupy
      ordinary \code{objN} handles (obj0--obj5 when the box is created
      first) and shift the indices of objects created after them;
      \code{BOX OFF} removes them and renumbers, \code{DEL} keeps the
      tracked wall indices renumbered, and \code{LIST} tags them
      \code{[wall: static, inverse\_mass=0]} so you can always tell
      which is which. Deleting a wall \textbf{dissolves} the box
      (\code{system.box} reads 0) but the surviving slabs stay tracked
      --- \code{LIST} keeps their \code{[wall]} tag, bare \code{BOX}
      reports the dissolved state, and \code{BOX <size>} /
      \code{BOX OFF} replaces / removes them without leaking an orphan
      slab.
\item \textbf{A failed \code{NEW} leaves nothing behind.} The
      initializer list is transactional: if any initializer --- or the
      final validation, e.g.\ a torus with inner $\ge$ outer --- fails,
      the half-built object is removed with the error (no ghost
      objects). Torus geometry in the initializer is resolved once at
      the closing brace, which is what makes the
      \code{inner\_radius}/\code{outer\_radius} pair genuinely
      order-independent.
\item \textbf{Magnetic torque is not potential-derived}: energy is not
      conserved while $M\!\cdot\!B$ acts (by design, matching legacy
      semantics).
\item \textbf{SPRK is translational-only} --- enforced by the gate.
\item \textbf{Static bodies} ($1/m = 0$) stop sourcing pairwise
      gravity, because their stored mass is 0 --- \code{BOX} walls
      included.
\item \textbf{\code{DEL} renumbers} subsequent object handles.
\item The terminal notebook has no cursor-key cell editing (pure
      \code{std}); JupyterLab provides that experience.
\item \textbf{Scene reverse is bounded}: the history ring keeps the
      last 20{,}000 forward frames; you cannot rewind past what was
      recorded (and \code{SCENE STOP} clears the ring).
\item \textbf{The scene evolves a \emph{copy}} of the notebook's
      system. Notebook \code{STEP}/\code{RUN} do not move the window,
      and scene playback does not move the notebook's state ---
      \code{SCENE REFRESH} re-syncs the window from the notebook
      whenever you want them aligned.
\item \textbf{Scene numeric arguments are term-level}:
      \code{scene rotate 15 -5} means ``yaw 15, pitch $-5$''; to pass a
      sum, parenthesize --- \code{scene zoom (1 + 0.5)}.
\item \textbf{Localhost only, one scene server per posim process} ---
      the window is a private local page, not a shared web service; a
      second \code{SCENE CREATE} reports the existing URL instead of
      opening a second server.
\item The window needs a browser with JavaScript enabled (any current
      Firefox/Chrome/Safari; nothing is fetched from the internet).
\end{itemize}

\end{document}
